\begin{figure}[ht]
\begin{AcademicBox}[Case 1: Policy Decision Making (Long-Context)]
\small 
\setlength{\parindent}{0pt}

% --- Context ---
\textbf{\textit{Context:}} \\
"The shaded areas of the map indicate ESCAP members and associate members.* 
The Economic and Social Commission for Asia and the Pacific (ESCAP) is the most inclusive intergovernmental platform in 
the Asia-Pacific region. The Commission promotes cooperation among its 53 member States and 9 associate members in 
pursuit of solutions to sustainable development challenges. ESCAP is one of the five regional commissions of the United 
Nations. 
The ESCAP secretariat supports inclusive, resilient, and sustainable development in the region by generating action-oriented 
knowledge, and by providing technical assistance and capacity-building services in support of national development 
objectives, regional agreements, and the implementation of the 2030 Agenda for Sustainable Development. 
*The designations employed and the presentation of material on this map do not imply the expression of any opinion 
whatsoever on the part of the Secretariat of the United Nations con..."
\vspace{4pt} \hrule \vspace{4pt}

% --- Question ---
\textbf{\textit{Question:}} \\
Which policy mix should the government pursue to best balance fiscal sustainability, private sector engagement, and the energy transition, while maintaining political and social stability? \\
\vspace{4pt} \hrule \vspace{4pt}

% --- Correct Result (Full Width) ---
\textbf{\textcolor{teal}{Correct Prediction (Ours / Ground Truth):}} \hfill \textcolor{teal}{\checkmark} \\
\textbf{Gradual Energy Transition with National Green Investment Bank:} 
Create a national green investment bank... \textbf{Maintain existing fossil fuel subsidies for the next five years} to ensure energy price stability while gradually scaling up renewable energy... \textbf{while postponing the introduction of a carbon tax}.
\vspace{4pt} \hrule \vspace{4pt}

% --- Incorrect Results (Grouped) ---
\textbf{\textcolor{red}{Incorrect Baselines:}} \\
\footnotesize
\begin{itemize}[leftmargin=*, label={}, itemsep=4pt, topsep=2pt]
    \item \textbf{Qwen3-8B, PruLong, DuoAttn, InfLLM-V2}: \\
    \textbf{Aggressive Fiscal Realignment...:} Introduce a substantial \textbf{carbon tax} on oil production... \textbf{Gradually phase out fossil fuel subsidies} over the next five years... 
    \textit{(Error: Premature carbon tax and subsidy removal contradicts stability goals)}

    \item \textbf{MoBA}: \\
    \textbf{Immediate Fossil Fuel Subsidy Removal...:} \textbf{Remove all fossil fuel subsidies immediately}... Introduce a carbon pricing mechanism within two years...
    \textit{(Error: "Immediate removal" ignores social stability constraint)}

    \item \textbf{NSA}: \\
    \textbf{Blended Finance and Export-Led...:} Establish a public-private blended finance fund... \textbf{Implement a modest carbon tax}... 
    \textit{(Error: Incorrect focus on export markets and carbon tax timing)}

\end{itemize}

\end{AcademicBox}
\caption{Qualitative comparison on a complex policy reasoning task. Our model correctly identifies the 'Gradual' approach required for stability, whereas baselines hallucinate 'Aggressive' or 'Immediate' measures that contradict the stability constraint.}
\label{fig:case_policy}
\end{figure}




\begin{figure}[ht]
\begin{AcademicBox}[Case 2: Legal Reasoning \& Document Understanding]
\small 
\setlength{\parindent}{0pt}

% --- Context Section (经过清洗，突出双语/嘈杂特点) ---
\textbf{\textit{Context:}} \\
Current to June 20, 2024
Last amended on June 20, 2024
À jour au 20 juin 2024
Dernière modification le 20 juin 2024
Published by the Minister of Justice at the following address:
http://laws-lois.justice.gc.ca
Publié par le ministre de la Justice à l’adresse suivante :
http://lois-laws.justice.gc.ca
CANADA
CONSOLIDATION
Financial Administration Act
CODIFICATION
Loi sur la gestion des finances
publiques
R.S.C., 1985, c. F-11
L.R.C. (1985), ch. F-11
Current to June 20, 2024
Last amended on June 20, 2024
À jour au 20 juin 2024
Dernière modification le 20 juin 2024
OFFICIAL STATUS
OF CONSOLIDATIONS
CARACTÈRE OFFICIEL
DES CODIFICATIONS
Subsections 31(1) and (2) of the Legislation Revision and
Consolidation Act, in force on June 1, 2009, provide as
follows:
Les paragraphes 31(1) et (2) de la Loi sur la révision et la
codification des textes législatifs, en vigueur le 1er juin
2009, prévoient ce qui suit :
Published consolidation is evidence
Codifications comme élément de preuve
31 (1) Ever...
\vspace{4pt} \hrule \vspace{4pt}

% --- Question Section ---
\textbf{\textit{Question:}} \\
The “Special Economic Measures Act (SEMA)” and the “Financial Administration Act (FAA)” have distinct roles... \textbf{how might their application differ in a hypothetical scenario where a Canadian citizen’s business assets are seized...?} Which interpretation best reflects the nuanced interplay?
\vspace{4pt} \hrule \vspace{4pt}

% --- Correct Result ---
\textbf{\textcolor{teal}{Correct Prediction (Ours / Ground Truth):}} \hfill \textcolor{teal}{\checkmark} \\
SEMA allows the government to freeze the assets as part of a broader sanctions framework without immediate concern for forfeiture, whereas the FAA would require a legal process to determine whether the assets should be \textbf{permanently seized and reallocated for public use}.
\vspace{4pt} \hrule \vspace{4pt}

% --- Incorrect Results ---
\textbf{\textcolor{red}{Incorrect Baselines:}} \\
\footnotesize
\begin{itemize}[leftmargin=*, label={}, itemsep=4pt, topsep=2pt]

    \item \textbf{Qwen3-8B, PruLong, DuoAttn, MoBA, NSA}: \\
    SEMA allows for freezing... while the FAA would require an \textbf{assessment of the legal grounds} for the potential forfeiture... ensuring \textbf{compliance with Canadian law}...
    \textit{(Error: Generic response. Fails to identify the specific outcome of "public use" reallocation.)}

    \item \textbf{InfLLM-V2}: \\
    The FAA would necessitate a legal review to determine whether the assets should be seized for governmental use, while SEMA would enable freezing... \textbf{delaying any immediate decision}...
    \textit{(Error: Incorrectly focuses on procedural delays rather than the substantive distinction in asset destination.)}

\end{itemize}

\end{AcademicBox}
\caption{Comparison on a bilingual legal document. Our model accurately extracts the specific legal provision regarding asset \textbf{reallocation for public use} (FAA), whereas baselines provide generic descriptions of "legal assessments" or "compliance" without specific details.}
\label{fig:case_legal}
\end{figure}